\documentclass[11pt,a4paper]{article}

\usepackage[margin=2.3cm]{geometry}
\usepackage{amsmath,amssymb}
\usepackage{graphicx}
\usepackage{booktabs}
\usepackage{array}
\usepackage{xcolor}
\usepackage{enumitem}
\usepackage[hidelinks]{hyperref}
\usepackage{caption}

\emergencystretch=1.5em
\captionsetup{font=small,labelfont=bf,width=.92\textwidth}

\setlength{\parskip}{4pt plus 1pt}
\setlength{\parindent}{0pt}

\newcommand{\code}[1]{\texttt{#1}}
\newcommand{\Gf}{G_{\mathrm{f}}}
\newcommand{\Gb}{G_{\mathrm{b}}}
\newcommand{\topk}{\operatorname{TopK}}
\newcommand{\rms}{\operatorname{RMS}}
\newcommand{\ok}{\checkmark}
\newcommand{\fail}{\texttimes}

\title{
Why a 24-layer stream-bottleneck model does not train\\[2pt]
\large Research journal
}
\author{}
\date{Started 2026-09-30 (box \code{lisbon}, 8$\times$RTX 5090).}

\begin{document}
\maketitle
\vspace{-2em}
\tableofcontents

% ============================================================
\section{The question}
% ============================================================

Every block of the models studied here ends in a stream bottleneck
(\code{placement=residual\_out}): the residual stream is replaced by
\begin{equation}
  B(h) = g_D\, W_D\, P_{S(h)}\, W_E\, h ,
  \qquad S(h) = \topk_K\bigl(|W_E h|\bigr),
  \label{eq:B}
\end{equation}
where \(W_E\in\mathbb{R}^{N\times d}\) is the encoder, \(W_D\in\mathbb{R}^{d\times N}\)
the decoder, \(P_S\) the coordinate projection onto the \(K\) kept features and
\(g_D\) a fixed scalar (\(\sqrt{d/K}\) under
\code{bottleneck\_decoder\_scale=backward\_preserving}, ``bp'').  Nothing routes
around \(B\): the stream entering block \(\ell+1\) is
\(x_{\ell+1}=\mathcal N\bigl(B(h_\ell)\bigr)\) with \(h_\ell = x_\ell+\Delta_\ell\),
\(\Delta_\ell\) the attention and MLP contributions of block \(\ell\), and
\(\mathcal N\) either the identity or an RMSNorm (\code{post\_norm}).
Throughout: \(d=1024\), \(N=4096\), \(K=512\) unless stated,
magnitude--direction decoupling (MD), TinyStories, lr \(3\cdot10^{-4}\),
cosine over 20k steps, batch \(64\times512\).

Observed before this study:

\begin{center}\small
\begin{tabular}{p{5.6cm}lp{6.4cm}}
\toprule
run & layers & outcome \\
\midrule
\code{al\_8L\_hard\_k512\_nopnorm\_bp\_std} & 8 & trains: val CE 3.46 (500), 2.34 (2k), 1.418 (20k) \\
\code{al\_500m\_hard\_k512\_pnorm\_bp\_std} & 24 & val CE 5.965 (500) \(\to\) 5.949 (10.5k), flat \\
\code{po\_500m\_hard\_k512\_pnorm\_bp\_std\_shared} & 24 & same plateau (5.954 at 2.5k) \\
\code{al\_500m\_hard\_k32\_pnorm\_md} & 24 & same plateau (5.947 at 4k) \\
\code{al\_500m\_hard\_k512\_nopnorm\_bp\_std} & 24 & CE \(=10.8249=\ln 50257\) from step \(\sim\)380, zero gradient \\
\code{al\_500m\_hard\_k512\_nopnorm\_bp\_ortho} & 24 & the same, from step \(\sim\)560 \\
\code{al\_500m\_hard\_k512\_nopnorm\_bp\_std\_lr6e4} & 24 & the same \\
\bottomrule
\end{tabular}
\end{center}

\textbf{The plateau is the unigram entropy.}  The token-frequency entropy of
\code{train.bin} is \(5.9665\) nats (and the cross-entropy of \code{val.bin}
under the train unigram is \(5.9665\)).  The post-norm 24-layer runs therefore
learn the token frequencies within \(\sim\)150 steps and nothing that depends on
the input afterwards.  The no-post-norm runs end at exactly \(\ln V\): uniform
logits.

% ============================================================
\section{Entry 1 --- the gains of one bottleneck}
% ============================================================

\textbf{Two exact properties of \eqref{eq:B}.}  Abs-\(\topk\) selection does not
change under a positive rescaling of \(h\), so \(B(\alpha h)=\alpha B(h)\); and
\(B\) is piecewise linear, \(B(h)=J(h)\,h\) with
\(J(h)=g_D W_D P_{S(h)} W_E\), which is also the Jacobian autograd uses (the mask
has zero derivative almost everywhere).

\textbf{Two gains.}  Per token, the forward (signal) gain and the backward gain
of a gradient \(g\) arriving at the output are
\begin{equation}
  \Gf=\frac{\|B(h)\|^2}{\|h\|^2},\qquad
  \Gb=\frac{\|J(h)^\top g\|^2}{\|g\|^2}.
\end{equation}
At initialization the entries of \(W_E,W_D\) are i.i.d.\ \(\mathcal N(0,1/d)\)
(projected to the MD sphere), so the pre-activations \(z=W_Eh\) are Gaussian with
variance \(\sigma^2=\|h\|^2/d\), and the decoder columns are orthonormal up to
\(O(d^{-1/2})\).  Then \(\|B(h)\|^2\approx g_D^2\sum_{i\in S}z_i^2\), while a
gradient that the selection did not see (isotropic, independent of \(S\)) puts
\(\|g\|^2/d\) on each kept coordinate:
\begin{equation}
  \Gf \approx g_D^2\,\frac{K}{d}\,s^2(\rho),\qquad
  \Gb \approx g_D^2\,\frac{K}{d},\qquad
  s^2(\rho)=\mathbb E\bigl[Z^2\,\big|\,|Z|>t_\rho\bigr]
  = 1+\frac{2\,t_\rho\,\varphi(t_\rho)}{\rho},
  \label{eq:s2}
\end{equation}
with \(\rho=K/N\), \(t_\rho=\Phi^{-1}(1-\rho/2)\) and \(\varphi\) the standard
normal density.  \(s^2\) is the \emph{selection gain}: the kept coefficients are
the largest ones, so their mean square exceeds that of a typical coefficient.
\(s^2 = 4.019\) at \(\rho=1/8\) and \(8.898\) at \(\rho=1/128\) (\(K=32\)).

\textbf{Check against the existing K sweep} (\code{docs/figures/diag500m/module\_gain},
24 layers at init, geometric mean over layers):

\begin{center}\small
\begin{tabular}{rccccc}
\toprule
\(K\) & \(s^2(K/N)\) & \(\Gf/\Gb\) (post-norm, \(g_D=1\)) & \(\Gf\) (bp, ortho) & \(\Gb\) (bp, ortho) \\
\midrule
32   & 8.898 & 9.159 & 8.869 & 2.156 \\
128  & 6.407 & 6.474 & 6.401 & 1.510 \\
512  & 4.019 & 4.033 & 4.018 & 1.006 \\
1024 & 2.894 & 2.926 & 2.896 & 0.999 \\
2048 & 1.857 & 1.871 & 1.858 & 1.000 \\
\bottomrule
\end{tabular}
\end{center}

With a post-norm the ratio \(\Gf/\Gb\) equals \(s^2(K/N)\) to 0.5--3\,\% at
every \(K\).  Without one (right columns) \(\Gf\) still equals \(s^2\) to
0.3\,\%; \(\Gb\) exceeds 1 at small \(K\) because there the gradient arriving
from the next bottleneck is no longer isotropic (it is concentrated on the
direction of the signal, which receives the forward gain).

\textbf{Invariance.}  Multiplying \(W_E\), \(W_D\) or \(g_D\) by scalars
multiplies \(\Gf\) and \(\Gb\) by the same factor; an RMSNorm after \(B\) divides
both by \(\Gf\); a tight-frame (orthogonal) initialization leaves the marginal of
\(z\) Gaussian; sharing \(W_E,W_D\) across layers changes neither.  The ratio
\(\Gf/\Gb=s^2\) is therefore unchanged by every intervention tried before this
study (post-norm, bp scale, MD gain spread \(\alpha\), orthogonal init, shared
projections).  Measured at init with \code{analysis/depth\_probe.py}: the
per-layer factor of the gradient profile is \(\times2.095\) (baseline),
\(\times2.096\) (shared), \(\times2.102\) (orthogonal).

% ============================================================
\section{Entry 2 --- consequences at depth}
% ============================================================

\textbf{Without a norm} the stream energy grows as \(\prod_\ell \Gf \approx s^{2L}\)
(with bp).  \(4.02^{24}\approx 3\cdot10^{14}\): measured stream RMS at the last
layer at init \(1.6\cdot10^{7}\) (predicted \(4.02^{12}=1.7\cdot10^7\)).  The final
RMSNorm divides the loss gradient by that RMS, so every parameter gradient is
\(10^{-11}\)--\(4\cdot10^{-8}\) (fc1 weight, layer 23 and layer 0), at or below
Adam's \(\epsilon=10^{-8}\).  Once training grows any gain, \(x^2\) exceeds the
float32 range inside the RMSNorm (\(|x|>1.8\cdot10^{19}\)), its output becomes
zero and so do the logits: CE \(=\ln V\) and zero gradient, as in the three
no-post-norm runs (their logged boundary gaps reach \(10^{18}\)--\(10^{24}\)).
At 8 layers the growth is \(4^{8/2}=256\) in RMS (measured 238) and training works.

\textbf{With a post-norm} the forward is normalized, and the backward through
each post-norm is divided by the same \(\sqrt{\Gf}\): per layer the gradient
energy is multiplied by \(\Gb/\Gf=s^{-2}\).  The gradient reaching block
\(\ell\) is \(s^{-(L-\ell)}\) times the last block's.  Measured at init (fc1
weight gradient RMS): \(6.7\cdot10^{-12}\) at layer 0 against
\(1.6\cdot10^{-4}\) at layer 23, a ratio \(2.4\cdot10^{7}\) (\(\times2.095\) per
layer; \(s=2.005\)).  At 8 layers the ratio is 247.

% ============================================================
\section{Entry 3 --- removing the selection gain (at initialization)}
% ============================================================

Two ways to make the forward stop carrying the selection bias, both now options
of \code{ActivationBottleneckConfig}:

\textbf{(a) Debiased values, \code{value\_shift}.}  Keep the support, shrink the
kept magnitudes,
\begin{equation}
  y_i=\operatorname{sign}(z_i)\,\max\bigl(|z_i|-\delta,0\bigr),\quad i\in S .
\end{equation}
The tangent Jacobian on the support is unchanged, so \(\Gb\) is unchanged, while
\(\Gf\) drops.  \code{energy} solves \(\delta\) per token from
\(\operatorname{mean}_S y^2=\operatorname{mean}_{\text{all}} z^2\), which gives
\(\Gf=\Gb\) exactly for a tight-frame encoder; \code{fixed} uses
\(\delta=\lambda^\ast\rms(z)\) with \(\lambda^\ast\) the root of
\(\mathbb E[(|Z|-\lambda)^2\mid |Z|>t_\rho]=1\) (\(\lambda^\ast=1.044\) at
\(\rho=1/8\)).

\textbf{(b) Code-space residual, \code{code\_residual}.}  Carry the \(K\)-sparse
code instead of re-encoding the decoded stream:
\begin{equation}
  c_0=\topk(W_E x_{\mathrm{emb}}),\qquad x_\ell=g_D W_D c_\ell,\qquad
  c_{\ell+1}=\topk\bigl(c_\ell+\alpha\,W_E\,\Delta_\ell(x_\ell)\bigr),
\end{equation}
one shared dictionary.  \(\topk\) is the Euclidean projection onto the set of
\(K\)-sparse vectors, so a block with \(\Delta_\ell=0\) is the identity on the
code, and \(\partial c_{\ell+1}/\partial c_\ell=P_{S_{\ell+1}}(I+\dots)\): the
residual connection of a pre-norm transformer, moved into code space.

Per-layer factor of the init gradient profile (fc1 weight gradient RMS,
\(\bigl(g_{23}/g_0\bigr)^{1/23}\)); 1 is flat:

\begin{center}\small
\begin{tabular}{lcc}
\toprule
variant & per-layer factor & last-layer stream RMS \\
\midrule
baseline, post-norm & 2.095 & 1.0 \\
baseline, no post-norm & 0.697 & \(1.6\cdot10^{7}\) \\
value shift \code{fixed} \(\lambda=0.7\), post-norm & 1.387 & 1.0 \\
value shift \code{fixed} \(\lambda^\ast=1.044\), post-norm & 1.044 & 1.0 \\
value shift \code{fixed} \(\lambda=1.3\), post-norm & 0.798 & 1.0 \\
value shift \code{energy}, post-norm & 1.044 & 1.0 \\
value shift \code{energy}, no post-norm & 0.939 & 6.4 \\
code residual \(\alpha=1\) / 0.3 / 0.1 & 1.008 / 0.995 / 0.997 & 8.7 / 3.0 / 2.1 \\
rblapsum surrogate \(T=0.3\) / 1 / 3, post-norm & 0.709 / 1.070 / 1.445 & 1.0 \\
lapsum surrogate \(T=1\), post-norm & 1.567 & 1.0 \\
\bottomrule
\end{tabular}
\end{center}

With the shift, the bottleneck's measured \(\Gf\) and \(\Gb\) are both
\(1.00\pm0.01\) at every layer.  The surrogate rows show the other side of the
same balance: a sharp surrogate (\(T=0.3\)) adds boundary-exchange gradient and
the backward explodes toward the input, a wide one vanishes; \(T\approx1\) is
close to critical for the backward only (its forward is the plain hard
\(\topk\)).

% ============================================================
\section{Entry 4 --- round 1: the post-norm variants still do not train}
% ============================================================

Eight 24-layer runs, the baseline's config otherwise, screened to 2000 steps
(\code{--stop-step}; the schedule stays defined over 20k steps).  Train CE:

\begin{center}\small
\begin{tabular}{lccccc}
\toprule
run & 100 & 200 & 300 & 500 & 1000 \\
\midrule
baseline (\code{al\_500m\_hard\_k512\_pnorm\_bp\_std}) & 6.43 & 5.98 & 5.97 & 5.98 & 5.99 \\
value shift energy, post-norm & 6.40 & 5.99 & 5.97 & 5.98 & \\
value shift fixed, post-norm & 6.40 & 5.98 & 5.97 & 5.99 & \\
value shift energy, post-norm, shared & 6.42 & 5.98 & 5.98 & 5.99 & \\
value shift energy, post-norm, \(K=32\) & 7.11 & 5.98 & 5.97 & & \\
rblapsum \(T=1\), post-norm & 6.41 & 5.98 & 5.97 & 5.98 & \\
value shift energy, \emph{no} post-norm & 6.53 & 5.41 & 4.85 & 4.25$^\dagger$ & \\
code residual \(\alpha=1\) & 5.25 & 3.79 & 3.28 & 2.79 & 2.19 \\
code residual \(\alpha=0.3\) & 5.13 & 3.75 & 3.29 & 2.84 & 2.21 \\
8 layers, no post-norm (reference) & 6.72 & 4.53 & 4.06 & 3.53 & 2.89 \\
\bottomrule
\end{tabular}\\[2pt]
\footnotesize$^\dagger$val CE at step 500.
\end{center}

\textbf{Measured:} every run with a post-norm follows the baseline to the second
decimal, including those whose initial gradient profile is flat (value shift)
or near-critical (rblapsum).  Both runs without a post-norm leave the plateau
from the start; the code residual is ahead of the 8-layer reference at every
step.  \textbf{Interpretation:} removing the selection gain is necessary (the
no-post-norm baseline overflows) but not sufficient; the post-norm itself stops
learning at this depth.

% ============================================================
\section{Entry 5 --- forensics of the failed run at step 2000}
% ============================================================

\code{analysis/depth\_probe.py --ckpt ... --init-reference} on
\code{al\_500m\_hard\_k512\_pnorm\_bp\_std/ckpt\_step2000.pt} and, for
contrast, \code{al\_8L\_hard\_k512\_nopnorm\_bp\_std/ckpt\_step2000.pt}.
\(\Delta W/W_0=\|W-W_0\|/\|W_0\|\) against the run's own step-0 weights
(rebuilt with the same seed; \code{md\_init\_} draws on the GPU, as in
\code{train()}).

\begin{center}\small
\begin{tabular}{lcccc}
\toprule
 & layers 0--12 & layer 18 & layer 21 & layer 23 \\
\midrule
\(\Delta W/W_0\), fc1 (24 layers, post-norm) & \(<5\cdot10^{-4}\) & 0.003 & 0.022 & 0.204 \\
fc1 gradient RMS & \(10^{-20}\)--\(10^{-16}\) & \(5\cdot10^{-15}\) & \(2\cdot10^{-12}\) & \(5\cdot10^{-8}\) \\
cosine between positions & 0.07--0.21 & 0.23 & 0.992 & 1.000 \\
bottleneck \(\Gf\) & 3.97--4.05 & 4.64 & 81 & 2680 \\
\midrule
\(\Delta W/W_0\), fc1 (8 layers, no post-norm) & \multicolumn{4}{c}{0.30--0.54 at every layer; cosine between positions \(\le0.17\)} \\
\bottomrule
\end{tabular}
\end{center}

\textbf{Measured:} after 2000 steps the first 13 blocks of the 24-layer run have
not moved (Adam with \(\epsilon=10^{-8}\) and gradients of \(10^{-20}\)--\(10^{-16}\));
the per-layer gradient factor has grown from 2.1 at init to 3.5.  In the last
three blocks the bottleneck's forward gain rises to 81--2680 (their outputs
carry 56--18861 times the energy of the unit carry, Entry 6), and the streams of
all positions are the same vector (cosine 1.000): the output cannot depend on the input, which is the
unigram plateau.

% ============================================================
\section{Entry 6 --- the second mechanism: a renormalized carry}
% ============================================================

\textbf{Common-mode fraction.}  For the streams \(v_t\) of the positions of a
sequence let \(f(v)=\|\bar v\|^2/\overline{\|v_t\|^2}\), \(\bar v\) the mean over
positions (\(t\ge16\)).  For \(v_t=c+u_t\) with independent \(u_t\) this is
\(\|c\|^2/(\|c\|^2+\|u\|^2)\), the expected cosine between two positions, and it
is what \code{depth\_probe.py} records as \code{f\_stream} (and, sampled, as
\code{tok\_cos}).  Per block it also records \(r=\overline{\|\Delta\|^2}/\overline{\|x\|^2}\)
and \(r_c=\|\bar\Delta\|^2/\overline{\|x\|^2}\), the energy of the block's
contribution and of its token-independent part relative to the carry.

\textbf{Model.}  Adding \(\Delta\) gives \(f(h)\approx(f(x)+r_c)/(1+r)\).  The
bottleneck maps a small common fraction by the slope of its correlation map;
for Gaussian codes and \(\phi(z)=z\,\mathbf 1[|z|>t]\) the Hermite expansion
gives
\begin{equation}
  \kappa=C_B'(0)=\frac{\mathbb E[Z\phi(Z)]^2}{\mathbb E[\phi(Z)^2]}
  =\mathbb E\bigl[Z^2\mathbf 1(|Z|>t_\rho)\bigr]=\rho\,s^2(\rho),
\end{equation}
the Gaussian energy fraction that survives the selection (0.50 at
\(\rho=1/8\), 0.070 at \(K=32\), 1 for \(K=N\)).  An RMSNorm leaves angles, and
so \(f\), unchanged, but it resets the carry to unit RMS, so \(r\) and \(r_c\)
stay \(O(1)\) at every depth and
\begin{equation}
  f_{\ell+1}=\kappa\,\frac{f_\ell+r_c}{1+r}
  \quad\Longrightarrow\quad
  f^\ast=\frac{\kappa\,r_c}{1+r-\kappa},
  \label{eq:fstar}
\end{equation}
approached geometrically.  Without a norm the carry accumulates energy, \(r_\ell\)
falls with depth, and there is no fixed point pulling \(f\) toward 1.

\textbf{Measured at initialization} (24 layers, \code{init3/}; \(r\), \(r_c\)
averaged over blocks 6--23, \(f\) over blocks 12--23):

\begin{center}\small
\begin{tabular}{lcccccc}
\toprule
 & \(\kappa=\rho s^2\) & \(\kappa\) measured & \(r_c\) & \(r\) & \(f^\ast\) \eqref{eq:fstar} & \(f\) measured \\
\midrule
post-norm, \(K=32\)   & 0.070 & 0.092 & 0.37 & 1.74 & 0.010 & 0.013 \\
post-norm, \(K=512\)  & 0.502 & 0.544 & 0.88 & 1.91 & 0.183 & 0.197 \\
post-norm, \(K=2048\) & 0.929 & 0.955 & 2.39 & 2.56 & 0.843 & 0.883 \\
post-norm, \(K=N\)    & 1     & 1.000 & 2.62 & 2.66 & 0.985 & 0.989 \\
\midrule
no norm, \(K=N\)            & 1 & 1.002 & \multicolumn{2}{c}{\(r\): 0.17 \(\to\) 0.05} & -- & 0.56 \\
dense, no bottleneck        & -- & 1.000 & \multicolumn{2}{c}{\(r\): 0.16 \(\to\) 0.05} & -- & 0.59 \\
value shift, no norm        & 0.502 & 0.520 & \multicolumn{2}{c}{\(r\): 0.15 \(\to\) 0.04} & -- & 0.009 \\
code residual               & -- & 0.883 & \multicolumn{2}{c}{\(r\): 0.06 \(\to\) 0.02} & -- & 0.06 \\
\bottomrule
\end{tabular}
\end{center}

\textbf{Measured:} with a post-norm, \eqref{eq:fstar} with the \emph{theoretical}
\(\kappa\) predicts the common-mode fraction to within 0.003--0.04 from \(K=32\)
to \(K=N\).  With \(K=N\) (a fixed random linear carry, no selection) the stream
is token-independent already at initialization (\(f=0.99\)); the same carry
without the post-norm stays at 0.56.  Most of \(r_c\) comes from the MLP (GELU
has a nonzero mean): the attention's share is 0.2 of 0.9 at \(K=512\).
\textbf{Interpretation:} the selection's decorrelation (\(\kappa<1\)) is what keeps
the \(K=512\) post-norm stack at \(f\approx0.2\) at initialization.  Training
makes the top blocks emit a large token-independent output (the unigram):
their \(r_c\) grows (to 56--18861 times the carry energy in the failed run,
table below), the selection becomes the same for every position (\(\kappa\to1\)), and
\(f\to1\) (cosine 1.000 at step 2000, Entry 5).  The large \(\Gf\) of those
blocks also divides the gradient reaching everything below them, which is why
the per-layer gradient factor grew from 2.1 to 3.5.

\textbf{The same quantities in the failed run} (\code{al\_500m\_hard\_k512\_pnorm\_bp\_std},
\(\kappa_\ell=f_{\ell+1}/f(h_\ell)\) measured):

\begin{center}\small
\begin{tabular}{lcccccccc}
\toprule
block & 0--17 & 18 & 19 & 20 & 21 & 22 & 23 \\
\midrule
\(f\) (step 2000)            & 0.07--0.23 & 0.23 & 0.36 & 0.80 & 0.992 & 1.000 & 1.000 \\
\(\kappa\)                  & 0.54--0.56 & 0.78 & 1.07 & 1.006 & 1.000 & 1.000 & -- \\
\(r_c\)                     & 0.47--0.90 & 1.19 & 3.04 & 15.3 & 56.4 & 528 & 18861 \\
\(r_c/r\)                   & 0.27--0.47 & 0.55 & 0.82 & 0.995 & 1.000 & 1.000 & 1.000 \\
selection energy \(s^2\)     & 4.02 & 4.06 & 4.53 & 6.24 & 7.79 & 8.00 & 7.98 \\
\bottomrule
\end{tabular}
\end{center}

Blocks 0--17 are still at their initialization values (they did not move).  The
blocks that trained emit an almost purely token-independent output
(\(r_c/r\to1\)); their inputs are then exactly \(K\)-sparse in the encoder
(\(s^2\to N/K=8\): all of the energy on the kept features), the support is the
same for every position (\(\kappa\to1\)), and \(f\to1\).  The recursion
\eqref{eq:fstar} reproduces it block by block (block 20: \((0.80+15.3)/16.4=0.98\),
measured 0.992).  Step 10000 is the same picture with larger top-block energies
(\(r_c=5\cdot10^5\) at block 23).

\textbf{The value shift under a post-norm collapses too} (\code{li\_24L\_k512\_pnorm\_shift\_energy},
step 2000, val CE 5.959; its initial gradient profile was flat, \(\times1.047\) per block):

\begin{center}\small
\begin{tabular}{lcccccccc}
\toprule
block & 0 & 6 & 12 & 16 & 18 & 20 & 22 & 23 \\
\midrule
\(f\)                 & 0.07 & 0.50 & 0.67 & 0.91 & 0.994 & 1.000 & 1.000 & 1.000 \\
\(\kappa\)            & 0.56 & 0.76 & 0.83 & 1.00 & 1.00 & 1.00 & 1.00 & -- \\
\(r_c/r\)             & 0.28 & 0.78 & 0.86 & 0.98 & 1.00 & 1.00 & 1.00 & 1.00 \\
bottleneck \(\Gf\)     & 1.00 & 1.07 & 1.11 & 1.53 & 2.23 & 5.82 & 244 & 470 \\
\(\Delta W/W_0\), fc1 & \(5.9\cdot10^{-3}\) & \(6.8\cdot10^{-3}\) & \(6.7\cdot10^{-3}\) & \(8.3\cdot10^{-3}\) & 0.015 & 0.020 & 0.063 & 0.19 \\
\bottomrule
\end{tabular}
\end{center}

\textbf{Measured:} with a balanced bottleneck every block moved in the first steps
(\(\Delta W/W_0\approx6\cdot10^{-3}\) at block 0, against \(4\cdot10^{-8}\) in the
baseline), but the token-independent share then spread down the stack --- from
block 16 on the streams of all positions coincide --- and the large forward gain
of the collapsed blocks (up to 470) brought the vanishing gradient back
(\(\times2.3\) per block at step 2000).  \textbf{Interpretation:} the renormalized
carry is sufficient on its own to stop learning at 24 layers; the selection gain
only decides how early it happens.

\textbf{Training controls} (\(K=N\), 24 layers): with the post-norm val CE
5.967 (500), 5.956 (1000), on the plateau; without it, \emph{[pending]}.

\textbf{All runs at step 2000} (\code{depth\_probe.py --ckpt} on each run's \code{ckpt\_step2000.pt}).
Per-block gradient factor \((g_{\mathrm{last}}/g_0)^{1/(L-1)}\) of the fc1
weight, relative weight change of block 0 and of the last block, cosine between positions at the last
block:

\begin{center}\small
\begin{tabular}{lcccc}
\toprule
run (24 layers unless marked) & grad.\ factor & \(\Delta W/W_0\), block 0 & \(\Delta W/W_0\), last & cosine, last \\
\midrule
TopK, post-norm (the failed run) & 3.48 & 4.4e{-}08 & 0.20 & 1.000 \\
value shift, post-norm & 2.30 & 5.9e{-}03 & 0.19 & 1.000 \\
\(K=N\), post-norm & 2.64 & 2.0e{-}03 & 0.16 & 1.000 \\
value shift, no norm & 0.93 & 3.1e{-}01 & 0.54 & 0.205 \\
code residual & 1.02 & 3.9e{-}01 & 0.39 & 0.120 \\
code residual, \(\alpha=0.3\) & 1.04 & 3.6e{-}01 & 0.45 & 0.111 \\
code residual, \(K=32\) & 1.04 & 3.8e{-}01 & 0.40 & 0.082 \\
code residual, 8 layers & 1.08 & 3.7e{-}01 & 0.41 & 0.106 \\
TopK, no norm, 8 layers & 1.04 & 3.0e{-}01 & 0.54 & 0.173 \\
\bottomrule
\end{tabular}
\end{center}

\textbf{Measured:} the post-norm decides the outcome.  Every run with a post-norm ends with the
positions collapsed (cosine 1.000) and a vanishing gradient (\(\times2.3\)--\(3.5\) per block, whatever its
initial profile); every run without one keeps the positions apart (\(\le0.2\)), a flat gradient
(\(\times0.93\)--\(1.08\)), and moves all of its blocks by 0.3--0.5.

% ============================================================
\section{Entry 7 --- what a stream bottleneck must do: the skip's job}
% ============================================================

A pre-norm transformer trains at depth because of two properties of its stream:
every block is the identity plus a small update (the input--output Jacobian of
the stack is close to an isometry), and nothing normalizes the carried stream
(only the branch inputs are normalized).  A stream bottleneck replaces the skip,
so it has to take over both roles.  The three requirements, and which variant
meets them:

\begin{center}\small
\begin{tabular}{lcccc}
\toprule
 & balanced (\(\Gf=\Gb\)) & carry not renormalized & identity when \(\Delta=0\) & 24 layers train \\
\midrule
TopK, post-norm (the failed runs)      & \fail\ (\(s^2\)) & \fail & \fail & \fail \\
TopK, no norm                          & \fail\ (\(s^2\)) & \ok   & \fail & \fail\ (overflow) \\
value shift, post-norm                 & \ok & \fail & \fail & \fail \\
\(K=N\) (linear carry), post-norm      & \ok & \fail & \fail & \fail \\
\(K=N\) (linear carry), no norm        & \ok & \ok   & \fail & \ok \\
value shift, no norm                   & \ok & \ok   & \fail & \ok\ (slow) \\
code residual                          & \ok & \ok   & \ok   & \ok \\
\bottomrule
\end{tabular}
\end{center}

The third column is idempotence: \(\topk\) is the Euclidean projection onto the
\(K\)-sparse set, \(\topk(c)=c\) for \(K\)-sparse \(c\), so carrying the code
makes a silent block the identity.  Re-encoding the decoded stream does not,
because \(W_EW_D\neq I\) (and cannot be on \(512\)-sparse codes in \(d=1024\)
with \(N=4096\): the leakage of \(W_E^\top W_D\) is comparable to the code at
\(K/d=1/2\)); every block is a new random rank-\(K\) map, which preserves
energy on average but spreads the singular values of the product over depth.
The value shift and the linear carry satisfy the first two requirements and
train; the code residual satisfies all three and trains fastest (Entry 8).

\textbf{The non-idempotent carry, measured} (\code{analysis/carry\_spectrum.py}:
exact Jacobian \(\partial x_L/\partial x_0\) of the carry alone, branches off,
fresh Gaussian bottlenecks at the MD init statistics, 3 draws).  Effective rank
\((\sum\sigma^2)^2/\sum\sigma^4\) out of \(d=1024\) and mean \(\sigma^2\):

\begin{center}\small
\begin{tabular}{lccccc}
\toprule
depth \(L\) & 1 & 4 & 8 & 16 & 24 \\
\midrule
value shift, no norm: effective rank & 256 & 78 & 41 & 21 & 14 \\
\phantom{value shift, no norm:} mean \(\sigma^2\) & 1.00 & 0.99 & 0.99 & 1.00 & 1.02 \\
TopK + post-norm: mean \(\sigma^2\) & 0.26 & \(3.9\cdot10^{-3}\) & \(1.5\cdot10^{-5}\) & \(2.1\cdot10^{-10}\) & \(2.9\cdot10^{-15}\) \\
TopK, no norm: effective rank & 255 & 8.1 & 1.01 & 1 & 1 \\
code residual (silent blocks) & \multicolumn{5}{c}{identity: all \(\sigma=1\), effective rank 1024} \\
\bottomrule
\end{tabular}
\end{center}

\textbf{Measured:} the balanced carry keeps the mean \(\sigma^2\) at 1, but its energy
concentrates on ever fewer directions (effective rank \(\approx 256\cdot L^{-0.9}\));
after 24 blocks the stack transmits the input along about 14 of 1024 directions.
The post-norm carry has the same spectrum shape scaled by \(s^{-2L}\); the plain
carry is a single spike along the signal (the selection gain).
\textbf{Interpretation:} this is the information the value-shift model has to
re-inject through attention and the MLPs at every block, and the likely reason it
trains more slowly than the code residual, whose carry is the identity.



% ============================================================
\section{Entry 8 --- round 2: controls, depth, and the value shift in training}
% ============================================================

\textbf{Sparsity-free controls} (\(K=N\), 24 layers, 2000 steps): with the post-norm val CE 5.959
(plateau, collapsed at step 2000: cosine 1.000, gradient \(\times2.64\) per block); without it 2.072.
The post-norm alone decides the outcome for the same carry.

\textbf{Depth sweep} (val CE at step 2000):

\begin{center}\small
\begin{tabular}{lcccc}
\toprule
layers & 8 & 12 & 16 & 24 \\
\midrule
TopK, post-norm & 1.862 & 5.958 & 5.959 & 5.959 \\
TopK, no norm   & 2.340 & 3.210 & 10.825 (\(\ln V\), guard stop 1240) & 10.825 (\(\ln V\)) \\
code residual   & 1.890 & 1.860 & 1.839 & 1.814 \\
\bottomrule
\end{tabular}
\end{center}

\textbf{Measured:} with the post-norm the stack trains at 8 layers --- better than without a norm --- and
sits on the plateau from 12 layers on, although at 12 layers the initial gradient ratio between the last
and the first block is only \(\sim2\cdot10^3\).  Without a norm the plain stack slows at 12 layers and
overflows at 16.  The code residual improves monotonically with depth.
\textbf{Interpretation:} the collapse, not the size of the initial gradient ratio, sets the post-norm
threshold: at 8 blocks the lower blocks learn token features before the top blocks' unigram output locks
the collapse in; from 12 blocks on they do not.

\textbf{The value shift drifts in training} (\code{li\_24L\_k512\_nopnorm\_shift\_energy}).  Val CE
2.859 (2k), 2.426 (3.5k), 2.137 (6k), with spikes at single evaluations (2.527 at 5k, 2.730 at 6.5k) and
raw gradient norms up to 15--23 (clipped at 1).  The fraction of kept features that the shift clamps to
zero rises steadily: 0.15 (1k), 0.26 (2k), 0.43 (3k), 0.51 (5k), 0.56 (6.9k), with \(\delta/\sigma\) from
1.04 to 1.95: the support is effectively \(\sim\)225 of \(K=512\).  The \code{fixed} rule
(\(\delta=\lambda^\ast\sigma\) throughout) stalls instead (val 3.764, 3.575, 3.619 at 1k, 1.5k, 2k), and a
shared dictionary does not help (3.443 at 1.5k against 3.205).
\textbf{Interpretation:} the shift removes the selection gain of \emph{Gaussian} codes.  Training makes the
kept magnitudes heavy-tailed, the selection gain of a trained code grows toward \(N/K=8\), and the
energy-matching shift must remove ever more of it (clamping), while a fixed shift falls out of balance.
The value shift therefore confirms the mechanism (it trains at 24 layers where plain TopK cannot) but
is not a robust fix; the code residual, which never reselects from a fresh encoding, has none of these
problems.

\textbf{The 20k runs} (val CE at step 20\,000): code residual 1.2658 (seed 2: 1.2668); code residual with an
encoder/decoder pair per block (\code{share\_projections=false}, 563M parameters, 1.06\,s/step) 1.2628, 1.797 at
2k; value shift without a norm 1.5063; code residual at \(K=32\) 1.517.  \textbf{Measured:} per-block projections
are ahead of the shared pair by 0.002--0.003 from step 17k on, three times the seed difference, for 55\,\% more
parameters.

% ============================================================
\section{Note on how the two fixes were arrived at}
% ============================================================

Both fixes were designed against the first mechanism (Entry 3), before the second was known.  The value shift
equalizes \(\Gf\) and \(\Gb\) directly.  The code residual came from the residual principle: a deep stack trains
when every block is the identity plus a small update, and the stream bottleneck breaks that because it re-encodes
the stream and reselects from fresh Gaussian coefficients in every block --- which is where the selection gain
comes from.  Carrying the code makes \(\topk\) idempotent on its input, so a silent block is the identity and there
is no fresh reselection.  It was run without a post-norm from the start, as a pre-norm residual stream would be.
The second mechanism was found afterwards, when the value shift failed with the post-norm and trained without
it (Entry 4).  In hindsight the recursion \eqref{eq:fstar} states what the second mechanism requires --- a carry
that is not renormalized, so that \(r_\ell\) falls with depth --- and both no-norm fixes meet it for the same
reason: neither normalizes the carry.  What the code residual adds is the identity carry (Entry 7).

% ============================================================
\section{Implementation notes}
% ============================================================

\begin{itemize}[leftmargin=1.2em]
\item \code{value\_shift} works on the \(K\) gathered values (only
  \(\sigma^2\) needs all \(N\)); gate forward+backward at \(8\times512\times4096\)
  bf16: 2.2\,ms against 1.0\,ms for plain \(\topk\) (7.9\,ms for the first,
  dense implementation, which the round-1 shift runs used).
\item A latent bug found on the way: \(g_D\) is a plain attribute that only
  \code{md\_init\_} set, so \code{load\_for\_inference} and the analysis
  \code{--ckpt} paths evaluated every bp checkpoint with \(g_D=1\) --- harmless
  behind a post-norm, a different network without one (e.g.\ the 8-layer
  reference).  The controller now restores it at install
  (\code{f9a869c}); training is unchanged (\code{tools/repro\_check.py}
  identical in all three modes).
\end{itemize}

\end{document}
